\documentclass[12pt]{article}
\usepackage[utf8]{inputenc}
\usepackage[T1]{fontenc}
\usepackage{hyperref}
\title{Sansu Language Resources for K-12 Curriculum Development}
\author{Compiled by Claude Sonnet 5.5 (claude-sonnet-5-5) for Andr\'as Kornai}
\date{}
\begin{document}
\maketitle

\section{Introduction}
The code sca (Sansu) is retired in ISO 639-3. Glottolog records it as merged into Hlersu [hle] (change request 2011-076, effective 2012-02-03), on the grounds that Sansu is another name for Hlersu and its speakers call themselves Hlersu. IANA lists sca as deprecated with preferred value hle. The ``Bookkeeping'' family in the prompt is Glottolog's placeholder for such retired codes. It is not a language family, and there is no separate ``Sansu'' language to build a curriculum for. Hlersu is a Loloish (Sino-Tibetan) language of Yunnan, China, with about 15,000 speakers according to the Endangered Languages Project, which also cites a UNESCO classification of Vulnerable. Sansu is the name used for the variety in Yuanjiang County. This dossier therefore points to the Hlersu (hle) record and does not claim any Sansu-specific resources.

\section{Primary Educational Resources}
None found.

\section{Linguistic and Reference Materials}
\begin{itemize}
\item Glottolog record for the retired code, merged into Hlersu: \url{https://glottolog.org/resource/languoid/iso/sca}
\item Glottolog languoid page, ``Bookkeeping: Sansu'': \url{https://glottolog.org/resource/languoid/id/sans1276}
\item IANA deprecated subtag record for sca: \url{https://www.iana.org/assignments/lang-subtags-templates/sca.txt}
\item Wikipedia, Hlersu language: \url{https://en.wikipedia.org/wiki/Hlersu_language}
\item Endangered Languages Project entry for Hlersu, which cites Xu Xianmin and Bai Bibo (2013, Ethnic Publishing House, Beijing) as the main documentation: \url{https://www.endangeredlanguages.com/lang/5125}
\item Numeral data for Hlersu (Max Planck): \url{https://lingweb.eva.mpg.de/channumerals/Hlersu.htm}
\end{itemize}

\section{Religious and Community Resources}
None found.

\section{Digital and Audio Resources}
None found.

\section{Curriculum Development Context}
Any curriculum work should be recorded under Hlersu (hle). This was not researched beyond the records above, so the situation of Hlersu in Yunnan schools is unknown.
\end{document}